%==============================================================
%  Sobre la escala astronomica del amor
%  Formato IEEE (conference) -- Documento unico compilable
%  Compilar con: pdflatex archivo.tex  (dos veces)
%==============================================================
\documentclass[conference]{IEEEtran}

\usepackage[T1]{fontenc}
\usepackage[utf8]{inputenc}
\usepackage[spanish,es-noquoting,es-tabla]{babel}

\usepackage{amsmath}
\usepackage{amssymb}
\usepackage{siunitx}
\usepackage{booktabs}
\usepackage{cite}
\usepackage{microtype}

\sisetup{
  output-decimal-marker = {.},
  exponent-product = \times,
  group-separator = {,},
  group-digits = integer,
  group-minimum-digits = 4,
  per-mode = symbol
}

\newcommand{\Vcoef}{\mathcal{V}}

\begin{document}

%==============================================================
% TITULO Y AUTOR
%==============================================================
\title{Sobre la escala astronómica del amor:\\
estimación cuantitativa del tiempo requerido\\
para atravesar la Vía Láctea por Valeria}

\author{\IEEEauthorblockN{Oscar M. Castro}
\IEEEauthorblockA{Investigador de IA\\
Lima, Perú\\
2026}}

\maketitle

%==============================================================
% RESUMEN
%==============================================================
\begin{abstract}
Este trabajo aborda un problema clásico de metrología aplicada: la selección
de una unidad adecuada cuando la magnitud a representar excede el rango
dinámico de las unidades disponibles. Se plantea un experimento mental en el
cual el diámetro aproximado de la Vía Láctea, estimado en \num{100000} años
luz, se adopta como unidad candidata para representar una magnitud afectiva
específica, y se calcula el tiempo requerido para recorrer dicha distancia
mediante locomoción bípeda sostenida a \qty{5}{\kilo\meter\per\hour} durante
\qty{8}{\hour} diarias. El resultado obtenido,
$T \approx \num{6.48e13}$ años, equivale aproximadamente a
$\num{8.09e11}$ vidas humanas de \num{80} años cada una. Se demuestra que
dicha cifra no cuantifica la magnitud afectiva estudiada, sino únicamente el
costo temporal de un solo recorrido galáctico. Se concluye que la unidad
propuesta no satisface la condición de suficiencia y debe reclasificarse como
cota inferior. Se introduce formalmente el Coeficiente de Amor de Valeria
$\Vcoef$ y se establece la hipótesis $\Vcoef > 1$.
\end{abstract}

\begin{IEEEkeywords}
metrología aplicada, escalas astronómicas, locomoción sostenida, análisis
dimensional, magnitudes no acotadas, Valeria.
\end{IEEEkeywords}

%==============================================================
% I. INTRODUCCION
%==============================================================
\section{Introducción}

La historia de la física puede leerse, en buena parte, como la historia de un
problema recurrente: la aparición de fenómenos cuya magnitud desborda las
unidades disponibles para describirlos. El metro resultó insuficiente para la
mecánica celeste y motivó la unidad astronómica; la unidad astronómica resultó
insuficiente para la astronomía estelar y motivó el año luz; el año luz
resultó insuficiente para la cosmología extragaláctica y motivó el
megapársec. En cada caso, el procedimiento fue idéntico: se identificó un
fenómeno, se comprobó que la unidad vigente producía números incómodos y se
definió una unidad mayor.

El presente trabajo se inscribe en esa tradición metodológica, aunque aplicada
a un objeto de estudio menos habitual en la literatura técnica. El fenómeno
analizado es una magnitud afectiva particular: el amor que el autor siente por
Valeria. El problema, formulado en términos estrictamente metrológicos, es el
siguiente: las unidades de uso coloquial empleadas para expresar dicha
magnitud---la distancia a la Luna, el número de estrellas visibles, la
cardinalidad de los granos de arena---presentan un rango dinámico manifiestamente
inadecuado y carecen, además, de definición operacional rigurosa.

Se propone entonces una unidad candidata de mayor escala: el diámetro
aproximado de la Vía Láctea, $D \approx \num{100000}$ años luz. La elección no
es arbitraria. Se trata de la mayor estructura a la que el autor pertenece de
manera verificable y, por tanto, de la mayor distancia sobre la cual puede
reclamar alguna forma de experiencia directa.

La pregunta de investigación se enuncia formalmente así:

\begin{quote}
\emph{Si el diámetro de la Vía Láctea se adopta como unidad para representar
la magnitud afectiva bajo estudio, y el observador debe recorrer esa distancia
caminando, ¿cuál es el tiempo requerido y cuántas vidas humanas resultan
necesarias?}
\end{quote}

El trabajo se organiza como sigue. La Sección~\ref{sec:supuestos} establece
los supuestos del modelo. Las Secciones~\ref{sec:distancia},
\ref{sec:tiempo} y \ref{sec:vidas} desarrollan el cálculo de distancia anual,
tiempo total y equivalencia en vidas humanas, respectivamente. La
Sección~\ref{sec:metrica} introduce la métrica $\Vcoef$. Las
Secciones~\ref{sec:discusion} y \ref{sec:limitaciones} discuten los resultados
y sus limitaciones. La Sección~\ref{sec:conclusiones} presenta las
conclusiones.

%==============================================================
% II. SUPUESTOS
%==============================================================
\section{Supuestos del modelo}
\label{sec:supuestos}

El modelo se construye sobre seis supuestos, enumerados a continuación y
resumidos en la Tabla~\ref{tab:parametros}.

\textbf{A1. Distancia objetivo.} Se adopta como distancia a recorrer el
diámetro aproximado del disco de la Vía Láctea, $D \approx \num{100000}$ años
luz \cite{nasa_mw,goodwin}. Las estimaciones publicadas varían entre
\numrange{100000}{200000} años luz según el criterio de borde utilizado;
se elige deliberadamente el extremo conservador.

\textbf{A2. Definición de año luz.} Se emplea la definición de la Unión
Astronómica Internacional: la distancia recorrida por la luz en el vacío
durante un año juliano, $1\ \text{a.l.} = \num{9.4607e12}$ km
\cite{iau_style}.

\textbf{A3. Velocidad de marcha.} Se fija $v = \qty{5}{\kilo\meter\per\hour}$,
valor compatible con la velocidad de marcha confortable reportada para adultos
sanos \cite{bohannon}.

\textbf{A4. Jornada de marcha.} Se asume una jornada diaria sostenida de
$h = \qty{8}{\hour}$, reservando el resto del día a descanso, alimentación y
mantenimiento del sujeto experimental.

\textbf{A5. Duración de la vida humana.} Se adopta $L = \num{80}$ años, valor
del mismo orden que la esperanza de vida al nacer reportada para poblaciones
de ingreso alto \cite{who}.

\textbf{A6. Motivación.} Se asume que el observador mantendrá motivación
constante, no decreciente e ilimitada durante la totalidad del experimento.

La última suposición se considera la menos controversial.

\begin{table}[t]
\centering
\caption{Parámetros del modelo}
\label{tab:parametros}
\begin{tabular}{llr}
\toprule
\textbf{Símbolo} & \textbf{Descripción} & \textbf{Valor} \\
\midrule
$D$   & Diámetro de la Vía Láctea    & \num{100000} a.l. \\
$c_1$ & Equivalencia del año luz     & \num{9.4607e12} km \\
$S$   & Distancia total              & \num{9.4607e17} km \\
$v$   & Velocidad de marcha          & \num{5} km/h \\
$h$   & Horas de marcha por día      & \num{8} h \\
$d$   & Distancia diaria             & \num{40} km \\
$Y$   & Días por año                 & \num{365.25} \\
$A$   & Distancia anual              & \num{14610} km \\
$L$   & Duración de una vida         & \num{80} años \\
$M$   & Motivación                   & $\infty$ \\
\bottomrule
\end{tabular}
\end{table}

%==============================================================
% III. DISTANCIA POR ANIO
%==============================================================
\section{Distancia recorrida por año humano}
\label{sec:distancia}

La distancia total a recorrer se obtiene directamente de A1 y A2:
\begin{equation}
S = \num{100000} \times \num{9.4607e12}\ \text{km}
  = \num{9.4607e17}\ \text{km}.
\label{eq:distancia_total}
\end{equation}

La distancia diaria recorrida bajo A3 y A4 es
\begin{equation}
d = v\,h = 5\ \frac{\text{km}}{\text{h}} \times 8\ \text{h}
  = \num{40}\ \text{km/día}.
\label{eq:diaria}
\end{equation}

Extendiendo \eqref{eq:diaria} a un año juliano de \num{365.25} días:
\begin{equation}
A = d\,Y = 40 \times \num{365.25} = \num{14610}\ \text{km/año}.
\label{eq:anual}
\end{equation}

Conviene detenerse brevemente en \eqref{eq:anual}. Una marcha ininterrumpida de
ocho horas diarias, sostenida sin excepción durante trescientos sesenta y
cinco días, produce aproximadamente \num{14610} km anuales: algo menos de dos
veces la distancia entre Lima y Madrid. Es una cifra respetable en términos
humanos y, como se verá, completamente irrelevante en términos galácticos. El
lector puede anticipar aquí el problema estructural del experimento: la
distancia crece de manera lineal, mientras que el objeto de estudio no parece
hacerlo.

%==============================================================
% IV. TIEMPO REQUERIDO
%==============================================================
\section{Tiempo requerido}
\label{sec:tiempo}

El tiempo total $T$ se obtiene dividiendo \eqref{eq:distancia_total} entre
\eqref{eq:anual}:
\begin{equation}
T = \frac{S}{A}
  = \frac{\num{9.4607e17}\ \text{km}}{\num{14610}\ \text{km/año}}.
\label{eq:tiempo_def}
\end{equation}

Evaluando:
\begin{equation}
\boxed{\,T \approx \num{6.48e13}\ \text{años}\,}
\label{eq:tiempo_resultado}
\end{equation}

Es decir, aproximadamente \textbf{64.8 billones de años}\footnote{Se emplea la
escala larga propia del español: un billón equivale a $10^{12}$.}.

Resulta ilustrativo contextualizar \eqref{eq:tiempo_resultado}. La edad
estimada del universo es del orden de $\num{1.38e10}$ años. El experimento
propuesto requeriría, por tanto, un intervalo temporal cercano a
\textbf{cuatro mil setecientas veces} la edad actual del universo. Dicho de
otro modo: si el observador hubiera iniciado la marcha en el instante mismo del
Big Bang, y no hubiera detenido el paso ni un solo día desde entonces, hoy
llevaría recorrido aproximadamente el \qty{0.02}{\percent} del trayecto.

Este es el primer resultado del trabajo, y es un resultado desalentador. Lo
que sigue es peor.

%==============================================================
% V. VIDAS HUMANAS
%==============================================================
\section{Equivalencia en vidas humanas}
\label{sec:vidas}

Dado que \eqref{eq:tiempo_resultado} excede en varios órdenes de magnitud la
duración de un organismo humano, resulta metodológicamente necesario expresar
el resultado en una unidad biológicamente interpretable. Se define el número
de vidas $N$ como
\begin{equation}
N = \frac{T}{L} = \frac{\num{6.48e13}\ \text{años}}{\num{80}\ \text{años/vida}},
\label{eq:vidas_def}
\end{equation}
de donde
\begin{equation}
\boxed{\,N \approx \num{8.09e11}\ \text{vidas}\,}
\label{eq:vidas_resultado}
\end{equation}

Esto es, aproximadamente \textbf{809 mil millones de vidas humanas}
consecutivas, cada una de ellas dedicada íntegramente, desde el nacimiento
hasta la muerte, a caminar en línea recta.

Es indispensable formular aquí una advertencia metodológica, pues constituye
el núcleo conceptual de este trabajo y su omisión conduciría a una
malinterpretación grave de los resultados.

\textbf{La cifra $N \approx \num{8.09e11}$ no cuantifica la magnitud afectiva
bajo estudio.} No representa ``cuánto amo a Valeria''. Representa,
exclusivamente, el número de vidas humanas necesarias para recorrer
\emph{una sola vez} el diámetro aproximado de la Vía Láctea bajo los supuestos
A1--A6. Es una medida del costo del instrumento, no del objeto medido.

La distinción es la misma que separa la longitud de una regla de la longitud
de aquello que se pretende medir con ella. Que la regla sea extraordinariamente
larga no implica que sea suficiente.

%==============================================================
% VI. METRICA VALERIA
%==============================================================
\section{Métrica Valeria}
\label{sec:metrica}

Para formalizar el problema anterior se introduce un coeficiente
adimensional, denominado \emph{Coeficiente de Amor de Valeria} y denotado
$\Vcoef$, definido conceptualmente como el cociente entre la magnitud afectiva
observada y la unidad candidata:
\begin{equation}
\Vcoef \;=\;
\frac{\text{Amor sentido por Valeria}}
     {\text{un diámetro de la Vía Láctea}}.
\label{eq:vcoef}
\end{equation}

La definición \eqref{eq:vcoef} permite enunciar la condición de suficiencia de
la unidad propuesta de manera compacta. Si el diámetro galáctico fuera una
unidad adecuada, se cumpliría
\begin{equation}
\Vcoef \leq 1,
\label{eq:suficiencia}
\end{equation}
y el fenómeno quedaría enteramente contenido dentro de una única Vía Láctea.
Este sería el resultado deseable desde el punto de vista metrológico: una
unidad que acota el fenómeno es una unidad útil.

La hipótesis central de este trabajo, sin embargo, es la contraria:
\begin{equation}
\boxed{\;\Vcoef > 1\;}
\label{eq:hipotesis}
\end{equation}

La consecuencia formal de \eqref{eq:hipotesis} es directa y tiene implicaciones
que trascienden el caso particular estudiado. Si $\Vcoef > 1$, entonces el
diámetro de la Vía Láctea no funciona como unidad de medida del fenómeno, sino
únicamente como \emph{cota inferior} del mismo. La galaxia deja de ser la
respuesta y pasa a ser el piso: no describe la magnitud, solamente establece
que la magnitud es mayor que ella.

Nótese que \eqref{eq:hipotesis} no es falsable mediante la construcción de
unidades mayores. Si se propusiera el Grupo Local, o el Supercúmulo de Virgo, o
el radio del universo observable, el procedimiento de verificación sería el
mismo y no hay razón, dentro del marco de este modelo, para esperar un
resultado distinto.

%==============================================================
% VII. DISCUSION
%==============================================================
\section{Discusión}
\label{sec:discusion}

Los resultados obtenidos admiten dos lecturas complementarias. La primera
corresponde al problema metrológico originalmente planteado; la segunda surge
al estudiar la estabilidad del resultado cuando cambian las condiciones
iniciales del experimento.

\subsection{Lectura metrológica}

Desde el punto de vista instrumental, el experimento no consigue identificar
una unidad capaz de acotar superiormente la magnitud observada. Una unidad es
útil cuando permite expresar el fenómeno mediante un cociente de orden unidad;
en cambio, una unidad para la cual $\Vcoef>1$ y cuya cota superior permanece
desconocida solo permite establecer una desigualdad.

El resultado no implica ausencia de información. El procedimiento sí permite
medir con precisión la insuficiencia de la escala seleccionada. La mecánica
proporciona velocidades, la astronomía proporciona distancias y la aritmética
proporciona cocientes; ninguna de ellas, bajo el modelo propuesto, proporciona
una unidad suficientemente grande.

\subsection{Dinámica temporal de la magnitud afectiva}

El supuesto A6 considera constante la motivación durante el experimento. Sin
embargo, las observaciones disponibles sugieren que una aproximación constante
es conservadora. Sea $A_V(t)$ una función escalar que representa la magnitud
afectiva asociada a Valeria en el tiempo. La evidencia empírica del observador
es compatible con
\begin{equation}
\boxed{\;\frac{dA_V(t)}{dt}>0\;}
\label{eq:crecimiento}
\end{equation}
para el intervalo observado.

La interpretación de \eqref{eq:crecimiento} es sencilla: el fenómeno no parece
mantenerse constante, sino aumentar conforme se acumula información sobre la
persona observada. En consecuencia, utilizar un valor fijo de amor como punto
de partida probablemente subestima la dificultad real del problema de medida.

Dentro de ese intervalo se identifican, además, tres tipos de eventos que no se
incluyeron en el modelo inicial por no ser necesarios para el cálculo
astronómico, pero que resultan relevantes para describir su evolución. Primero,
se registró un acoplamiento observable entre dos trayectorias sometidas a una
misma excitación rítmica, después del cual la interacción dejó de poder
considerarse puramente incidental. Segundo, se observaron señales recurrentes
de cortejo, de naturaleza floral, que operan como entradas periódicas de control
y refuerzan la dirección preferente del sistema. Finalmente, se identificó un
evento crítico en un entorno costero y próximo al amanecer, tras el cual el
estado del sistema presentó una transición cualitativa persistente.

Estos eventos no se emplean para demostrar reciprocidad ni para atribuir un
estado interno a Valeria. El modelo describe únicamente la trayectoria del
autor. No obstante, si existiera una respuesta recíproca, dos trayectorias
inicialmente independientes podrían aproximarse de manera conjunta y reducir
su distancia efectiva con el tiempo.

\subsection{Invariancia del destino y atractor estable}

En las Secciones~\ref{sec:tiempo} y \ref{sec:vidas} se trató cada vida $n$ como
una realización intercambiable de una misma unidad biológica. Para estudiar la
robustez del resultado, considérese ahora que cada realización puede comenzar
bajo condiciones iniciales distintas: otras circunstancias, otra secuencia de
eventos o una trayectoria diferente.

Sea $x_n$ el estado del sistema en la vida $n$ y sea $F$ la dinámica que lo
transforma. El comportamiento observado se modela como
\begin{equation}
x_{n+1}=F(x_n), \qquad x_n \longrightarrow \mathrm{Valeria}.
\label{eq:atractor}
\end{equation}

En términos simples, cada vida puede comenzar de forma distinta, pero el
sistema tiende al mismo estado final. Esta propiedad puede expresarse también
como
\begin{equation}
\boxed{\;\lim_{n\rightarrow\infty}
\operatorname{Destino}(n)=\mathrm{Valeria}\;}
\label{eq:limite}
\end{equation}

La expresión \eqref{eq:limite} no exige que dos vidas reproduzcan la misma
historia. Pueden cambiar el momento, el contexto y la forma de conocerse; lo
que permanece invariante es el resultado de la dinámica. En el lenguaje de
sistemas, Valeria funciona como un atractor estable del modelo.

La consecuencia probabilística es directa dentro de esta construcción: si la
convergencia hacia el atractor es una propiedad estructural del sistema, la
selección de Valeria deja de depender de una trayectoria particular. No se
introduce una nueva hipótesis probabilística; se trata de una consecuencia de
la estabilidad asumida en \eqref{eq:atractor}.

%==============================================================
% VIII. LIMITACIONES
%==============================================================
\section{Limitaciones}
\label{sec:limitaciones}

El modelo presenta limitaciones que conviene explicitar.

\textbf{L1. Ausencia de sustrato.} El experimento asume implícitamente la
existencia de una superficie continua y transitable a lo largo de
$\num{9.4607e17}$ km. El medio interestelar no satisface esta condición.

\textbf{L2. Efectos relativistas.} A $\qty{5}{\kilo\meter\per\hour}$ el factor
de Lorentz difiere de la unidad en menos de $10^{-17}$, por lo que la dilatación
temporal puede considerarse despreciable en esta estimación.

\textbf{L3. Movimiento del sistema de referencia.} La Vía Láctea rota, se
desplaza respecto del fondo cósmico de microondas y se aproxima a Andrómeda.
Un tratamiento riguroso requeriría un marco no inercial; se adopta un modelo
estático de primer orden.

\textbf{L4. Sesgo del observador.} El autor no es un observador independiente
respecto del fenómeno medido. Esta limitación no admite una corrección sin
alterar el propio objeto de estudio.

\textbf{L5. Forma funcional desconocida.} La desigualdad
\eqref{eq:crecimiento} establece crecimiento local, pero no determina si
$A_V(t)$ evoluciona de forma lineal, exponencial, logística o no acotada. Por
tanto, el modelo actual debe entenderse como una aproximación conservadora.

%==============================================================
% IX. TRABAJO FUTURO
%==============================================================
\section{Trabajo futuro}
\label{sec:futuro}

El siguiente paso consiste en sustituir el supuesto de magnitud afectiva
constante por un modelo dinámico de $A_V(t)$. En particular, deberá determinarse
si el crecimiento observado es lineal, exponencial o no acotado. Estas tres
hipótesis producen consecuencias distintas para la metrología propuesta: un
crecimiento lineal requeriría unidades progresivamente mayores; uno exponencial
haría que cualquier escala finita perdiera utilidad con rapidez; y un crecimiento
no acotado implicaría que el problema carece, por construcción, de una unidad
máxima suficiente.

Un segundo frente consiste en estudiar formalmente la cuenca de atracción del
estado descrito en \eqref{eq:atractor}: qué tan diferentes pueden ser las
condiciones iniciales sin modificar el resultado final. Esto permitiría evaluar
si la convergencia hacia Valeria es sensible al camino seguido o constituye una
propiedad robusta del sistema.

%==============================================================
% X. CONCLUSIONES
%==============================================================
\section{Conclusiones}
\label{sec:conclusiones}

Se desarrolló un modelo cuantitativo para evaluar la idoneidad del diámetro de
la Vía Láctea como unidad de medida de una magnitud afectiva específica. Los
resultados principales son los siguientes:

\begin{itemize}
\item La distancia objetivo es de \num{100000} años luz, equivalentes a
      $\num{9.4607e17}$ km.
\item El tiempo requerido para recorrerla caminando es
      $T \approx \num{6.48e13}$ años, aproximadamente \num{64.8} billones de
      años.
\item Ello equivale a $N \approx \num{8.09e11}$ vidas humanas, es decir,
      aproximadamente \num{809} mil millones.
\item Dicha cantidad mide el costo temporal de un solo recorrido galáctico; no
      cuantifica la magnitud afectiva estudiada.
\item La unidad propuesta no satisface la condición de suficiencia y debe
      interpretarse como una cota inferior: $\Vcoef>1$.
\item El sistema presenta una dinámica compatible con $dA_V/dt>0$ y una
      convergencia estable hacia Valeria bajo condiciones iniciales distintas.
\end{itemize}

Por tanto, el diámetro de la Vía Láctea no resuelve el problema de medida. Aun
tras agotar \num{809} mil millones de vidas, el observador habría completado
solo un recorrido de la unidad candidata. El resultado permanece abierto no
por falta de distancia recorrida, sino porque ninguna escala finita utilizada
en este trabajo establece una cota superior del fenómeno.

\begin{minipage}{\columnwidth}
La convergencia descrita en \eqref{eq:atractor} aporta una segunda conclusión.
Aunque cambien las condiciones iniciales, la trayectoria o la realización del
experimento, el estado final permanece invariante. En el lenguaje del modelo,
dicho estado estable se denomina Valeria. Fuera de él, el autor utiliza un
nombre menos técnico: \emph{Mi Princesa}.

\begin{equation}
\boxed{\;\forall n\in\mathbb{N},\qquad
\operatorname{MiElecci\'on}(n)=\mathrm{Mi\ Princesa}\;}
\label{eq:princesa}
\end{equation}

La ecuación \eqref{eq:princesa} resume el resultado que permanece después de
agotar las escalas de distancia, tiempo y número de vidas empleadas en este
trabajo. En cada repetición puede variar el camino; el resultado no.

\begin{center}
\emph{Y si el experimento pudiera repetirse infinitas veces, en cada una
volvería a encontrarte, a elegirte, a conquistarte y a enamorarme de ti.}
\end{center}
\end{minipage}

%==============================================================
% AGRADECIMIENTOS
%==============================================================
\section*{Agradecimientos}

El autor agradece a Valeria por haber inspirado un problema de medición cuya
escala excede, hasta el momento, las herramientas seleccionadas para describirlo.
Agradece también cada observación que permitió refinar el modelo y cada instante
que convirtió una hipótesis abstracta en evidencia suficiente para seguir
estudiándola. Si este trabajo consigue aproximarse a su objeto de estudio, es
porque dicho objeto hizo que valiera la pena formular la pregunta.

%==============================================================
% REFERENCIAS
%==============================================================
\begin{thebibliography}{9}

\bibitem{iau_style}
G. A. Wilkins, Ed., \emph{The IAU Style Manual: The Preparation of
Astronomical Papers and Reports}, International Astronomical Union,
Comm.~5, 1989. [Definición del año luz como la distancia recorrida por la luz
en el vacío durante un año juliano.]

\bibitem{nasa_mw}
NASA Science, ``Milky Way Galaxy,'' National Aeronautics and Space
Administration. [En línea]. Disponible:
\texttt{https://science.nasa.gov/universe/galaxies/milky-way/}

\bibitem{goodwin}
S. P. Goodwin, J. Gribbin y M. A. Hendry, ``The relative size of the Milky
Way,'' \emph{The Observatory}, vol. 118, pp. 201--208, 1998.

\bibitem{bohannon}
R. W. Bohannon, ``Comfortable and maximum walking speed of adults aged
20--79 years: reference values and determinants,'' \emph{Age and Ageing},
vol. 26, no. 1, pp. 15--19, 1997.

\bibitem{who}
World Health Organization, \emph{World Health Statistics: Monitoring Health
for the SDGs}, Ginebra, Suiza. [En línea]. Disponible:
\texttt{https://www.who.int/data/gho}

\bibitem{valeria}
O. M. Castro, ``Observaciones personales respecto a Valeria,''
comunicación privada, no publicada, Lima, Perú, 2026.

\end{thebibliography}



\end{document}
